\chapter{Współbieżność i asynchroniczność}
\label{chap:concurrency}

\section{Motywacja i podstawowe pojęcia}

W programie jednowątkowym instrukcje wykonywane są w jednym \emph{wątku sterowania}.
Nie znaczy to, że procesor nie wykonuje w tle innych czynności; chodzi o model
samego programu. W programie współbieżnym istnieje kilka niezależnych ciągów
instrukcji, których wykonanie może się przeplatać. Zamiast jednego zadania
rozpatrujemy np. odczyt danych, obliczenia i zapisywanie wyników.

\begin{description}[style=nextline]
\item[Współbieżność (\emph{concurrency})]
Wiele zadań może robić postępy w nakładających się okresach czasu.
Nie wymaga to, aby fizycznie wykonywały się w tej samej chwili.
\item[Równoległość (\emph{parallelism})]
Co najmniej dwa zadania wykonują się fizycznie jednocześnie, na przykład
na dwóch rdzeniach procesora.
\item[Asynchroniczność (\emph{asynchrony})]
Rozpoczęcie operacji nie wymusza natychmiastowego oczekiwania na jej wynik.
Wywołujący może kontynuować pracę, a wynik odebrać później.
\end{description}

Pojęcia te nie są synonimami. Można mieć współbieżność bez równoległości
(jeden rdzeń i przełączanie zadań), a także interfejs asynchroniczny bez nowego
wątku (np. opóźnione wykonanie funkcji).

\nota{Zegar i kolejność komunikatów w przykładach współbieżnych nie stanowią
gwarantowanego wyniku programu. Planista systemu może przeplatać wątki inaczej
przy każdym uruchomieniu.}

\subsection{Proces, wątek, zadanie i współdzielone dane}

Proces zwykle posiada własną przestrzeń adresową. Wątki tego samego procesu
współdzielą między innymi pamięć dynamiczną i zmienne globalne, ale mają własne
stosy i własne wykonanie instrukcji. \emph{Zadanie} to jednostka pracy,
którą można przypisać wątkowi; nie każde zadanie musi oznaczać nowy wątek.

Dostęp do wspólnego obiektu wymaga szczególnej ostrożności, jeżeli co najmniej
jeden wątek go modyfikuje. Sam fakt istnienia wielu wątków nie jest błędem;
problem pojawia się przy niezsynchronizowanym dostępie do pamięci.

\section{Uruchamianie wątków: \texttt{std::thread}}
\label{sec:thread}

Nagłówek \code{<thread>} udostępnia typ \code{std::thread} (C++11).
Utworzenie obiektu z funkcją uruchamia nowy wątek. Wywołanie \code{join()}
czeka na jego zakończenie; \code{joinable()} informuje, czy obiekt nadal
reprezentuje wątek wymagający obsłużenia.

\begin{lstlisting}[caption={Pierwszy program z dwoma wątkami (C++11).}]
#include <iostream>
#include <thread>

void work(int id) {
    std::cout << "Worker " << id << "\n";
}

int main() {
    std::thread a(work, 1);
    std::thread b(work, 2);

    a.join();
    b.join();
}
\end{lstlisting}

Kolejność wypisywania nie jest określona, a fragmenty wypisywane przez
różne wątki mogą się przeplatać. Jeżeli wymagana jest spójna linia wyjścia,
można chronić wypisywanie mutexem albo użyć \code{std::osyncstream} (C++20).

\subsection{Przekazywanie argumentów i czasu życia}

Argumenty przekazywane do konstruktora \code{std::thread} są odpowiednio
przechowywane na potrzeby nowego wątku; domyślnie nie oznacza to przekazania
modyfikowalnego argumentu przez referencję. Referencję można przekazać
za pomocą \code{std::ref} albo przechwycenia referencji w lambdzie.

\begin{lstlisting}[caption={Bezpieczne przekazanie referencji z zachowaniem jej czasu życia.}]
#include <functional>
#include <iostream>
#include <thread>

void increment(int& value) { ++value; }

int main() {
    int counter = 0;
    std::thread worker(increment, std::ref(counter));
    worker.join(); // counter nadal istnieje
    std::cout << counter << '\n'; // 1
}
\end{lstlisting}

Destruktor obiektu \code{std::thread}, który nadal jest
\code{joinable()}, wywołuje \code{std::terminate()}.
Dlatego przy wielu ścieżkach wyjścia lepiej stosować RAII,
np. \code{std::jthread}. Operacja \code{detach()} wymaga z kolei
gwarancji, że wątek nie będzie używał już zniszczonych obiektów.

\section{Wyścig danych, warunek wyścigu i relacja \emph{happens-before}}
\label{sec:race}

\term{Wyścig danych (\emph{data race})} występuje, gdy co najmniej dwa
potencjalnie współbieżne dostępy do tego samego miejsca pamięci są konfliktowe
(co najmniej jeden modyfikuje), nie są odpowiednio uporządkowane przez
synchronizację i nie są to dostępy atomowe w wymaganym sensie. W standardowym
C++ prowadzi to do \textbf{niezdefiniowanego zachowania}.

\begin{lstlisting}[caption={Błędny przykład: niezdefiniowane zachowanie.}]
// Fragment celowo niepoprawny; NIE uruchamiac jako poprawnego przykladu.
int counter = 0;
// Watki rownoczesnie wykonuja ++counter bez synchronizacji.
\end{lstlisting}

\term{Warunek wyścigu (\emph{race condition})} jest pojęciem szerszym:
wynik programu zależy od względnej kolejności zdarzeń. Warunek wyścigu nie
zawsze oznacza niezdefiniowane zachowanie; np. dwa poprawnie zsynchronizowane
wątki mogą rywalizować o to, który pierwszy wykona zadanie.

\term{Relacja \emph{happens-before}} formalizuje, jakie skutki wcześniejszych
operacji muszą być widoczne dla późniejszych. Tworzy się m.in. przez porządek
operacji w jednym wątku i przez operacje synchronizujące (np. odpowiednio
sparowane \code{release}/\code{acquire}, odblokowanie/zablokowanie mutexa).
To centralny element modelu pamięci C++.

\section{Ochrona danych: \texttt{mutex} i blokady RAII}
\label{sec:mutex}

\code{std::mutex} pozwala maksymalnie jednemu wątkowi utrzymywać blokadę
w danej chwili. Jeśli kilka operacji ma tworzyć niepodzielną \emph{sekcję
krytyczną}, umieszczamy je pod tą samą blokadą.

\begin{lstlisting}[caption={Bezpieczny licznik przy użyciu mutexa (C++11).}]
#include <iostream>
#include <mutex>
#include <thread>
#include <vector>

int main() {
    int counter = 0;
    std::mutex mtx;
    std::vector<std::thread> workers;

    for (int i = 0; i < 4; ++i) {
        workers.emplace_back([&] {
            for (int j = 0; j < 10000; ++j) {
                std::lock_guard<std::mutex> lock(mtx);
                ++counter;
            }
        });
    }
    for (auto& t : workers) t.join();
    std::cout << counter << '\n'; // 40000
}
\end{lstlisting}

\code{lock_guard} automatycznie przejmuje mutex w konstruktorze i zwalnia
w destruktorze. Dzięki temu zwolnienie blokady następuje także podczas
rozwijania stosu po wyjątku. \code{unique_lock} jest bardziej elastyczny:
pozwala jawnie odblokować i ponownie zablokować mutex oraz współpracuje
z \code{condition_variable}.

\subsection{Kilka mutexów i zakleszczenie}

Jeżeli dwa wątki blokują mutexy w przeciwnej kolejności, może wystąpić
\emph{deadlock}: pierwszy czeka na drugi mutex, drugi czeka na pierwszy.
W C++17 można bezpiecznie przejmować kilka mutexów przez \code{scoped_lock}.

\begin{lstlisting}[caption={Łączne przejęcie dwóch mutexów (C++17).}]
#include <mutex>

int main() {
    std::mutex a, b;
    {
        std::scoped_lock lock(a, b);
        // Operacje chronione jednoczesnie przez a i b.
    }
}
\end{lstlisting}

\nota{Mutexy nie chronią ,,samej zmiennej'' automatycznie: poprawność wymaga,
aby \emph{wszystkie} konfliktowe dostępy do chronionego stanu stosowały
uzgodniony mechanizm synchronizacji.}

\section{Oczekiwanie na stan: \texttt{condition\_variable}}
\label{sec:cv}

Mutex zapewnia wyłączność, ale nie sygnalizuje gotowości danych.
\code{std::condition_variable} pozwala czekać bez aktywnego odpytywania.
Do oczekiwania z predykatem używamy blokady
\begin{center}
\code{std::unique_lock<std::mutex>}.
\end{center}
Predykat opisuje stan, na który czeka konsument.

\begin{lstlisting}[caption={Producent i konsument z warunkiem gotowości (C++11).}]
#include <condition_variable>
#include <iostream>
#include <mutex>
#include <thread>

int main() {
    std::mutex mtx;
    std::condition_variable cv;
    bool ready = false;
    int data = 0;

    std::thread consumer([&] {
        std::unique_lock<std::mutex> lock(mtx);
        cv.wait(lock, [&] { return ready; });
        std::cout << data << '\n'; // 42
    });

    std::thread producer([&] {
        {
            std::lock_guard<std::mutex> lock(mtx);
            data = 42;
            ready = true;
        }
        cv.notify_one();
    });

    producer.join();
    consumer.join();
}
\end{lstlisting}

Wywołanie \code{wait(lock, pred)} ma następujący sens:
\begin{enumerate}
\item podczas trzymania mutexa sprawdzany jest predykat;
\item jeśli jest fałszywy, wątek atomowo zwalnia mutex i przechodzi w oczekiwanie;
\item po powiadomieniu (lub pozornym wybudzeniu) odzyskuje mutex i znów sprawdza predykat;
\item kończy czekanie dopiero wtedy, gdy predykat jest prawdziwy.
\end{enumerate}

Dopuszczalne są \emph{spurious wakeups}, czyli wybudzenia bez zmiany interesującego
nas stanu. Dlatego samo \code{cv.wait(lock)} zwykle nie wystarcza.
Powiadomienie \code{notify_one()} budzi co najwyżej jeden oczekujący wątek,
a \code{notify_all()} umożliwia wybudzenie wszystkich. Powiadomienia nie są
kolejkowane jak komunikaty; stan \code{ready}, a nie samo powiadomienie,
stanowi źródło prawdy.

\subsection{Kolejka producent--konsument}

W ogólniejszym przypadku producent umieszcza wiele elementów w kolejce,
a konsument czeka na warunek ,,kolejka niepusta lub zakończono produkcję''.
Flaga zakończenia jest konieczna, aby konsument nie oczekiwał bez końca
po przetworzeniu ostatniego elementu.

\begin{lstlisting}[caption={Współdzielona kolejka i sygnał zakończenia (C++11).}]
#include <condition_variable>
#include <iostream>
#include <mutex>
#include <queue>
#include <thread>

int main() {
    std::mutex m;
    std::condition_variable cv;
    std::queue<int> q;
    bool done = false;

    std::thread consumer([&] {
        for (;;) {
            int item;
            {
                std::unique_lock<std::mutex> lock(m);
                cv.wait(lock, [&] { return done || !q.empty(); });
                if (q.empty() && done) break;
                item = q.front();
                q.pop();
            }
            std::cout << "Consumed " << item << '\n';
        }
    });

    std::thread producer([&] {
        for (int x = 1; x <= 3; ++x) {
            {
                std::lock_guard<std::mutex> lock(m);
                q.push(x);
            }
            cv.notify_one();
        }
        {
            std::lock_guard<std::mutex> lock(m);
            done = true;
        }
        cv.notify_all();
    });

    producer.join();
    consumer.join();
}
\end{lstlisting}

\section{Operacje atomowe: \texttt{std::atomic}}
\label{sec:atomic}

\code{std::atomic<T>} udostępnia atomowe operacje na obiekcie \code{T}.
Nie oznacza to ogólnie ,,braku jakichkolwiek blokad'' w implementacji:
nie każda instancja atomika jest \emph{lock-free}. Można to sprawdzić
metodą \code{is_lock_free()}.

\begin{lstlisting}[caption={Licznik atomowy z porządkiem relaxed (C++11).}]
#include <atomic>
#include <iostream>
#include <thread>
#include <vector>

int main() {
    std::atomic<int> counter{0};
    std::vector<std::thread> threads;
    for (int i = 0; i < 4; ++i) {
        threads.emplace_back([&] {
            for (int j = 0; j < 10000; ++j)
                counter.fetch_add(1, std::memory_order_relaxed);
        });
    }
    for (auto& t : threads) t.join();
    std::cout << counter.load() << '\n'; // 40000
}
\end{lstlisting}

W przypadku samego licznika potrzebujemy niepodzielności inkrementacji,
a nie synchronizacji dodatkowych danych. Stąd wystarcza
\code{memory_order_relaxed}. Domyślnie operacje atomowe stosują
\code{memory_order_seq_cst}, który narzuca silniejsze gwarancje porządkowania.

\subsection{Operacja porównaj-i-zamień}

\code{compare_exchange_weak(expected, desired)} atomowo sprawdza, czy
obecna wartość jest równa \code{expected}. Przy powodzeniu zapisuje
\code{desired}; przy niepowodzeniu aktualizuje \code{expected} do wartości
odczytanej. Wersja \code{weak} może zgłosić niepowodzenie także pozornie,
dlatego często umieszcza się ją w pętli.

\begin{lstlisting}[caption={Atomowe podnoszenie wartości do zadanego minimum (C++11).}]
#include <atomic>
#include <iostream>

int main() {
    std::atomic<int> maximum{3};
    int desired = 10;
    int expected = maximum.load(std::memory_order_relaxed);

    while (expected < desired &&
           !maximum.compare_exchange_weak(
               expected, desired,
               std::memory_order_relaxed)) {
        // expected zawiera najnowsza odczytana wartosc
    }
    std::cout << maximum.load() << '\n'; // 10
}
\end{lstlisting}

\subsection{Porządek pamięci: \texttt{release} i \texttt{acquire}}

Atomowość jednej zmiennej nie wystarcza, gdy jednym sygnałem chcemy
bezpiecznie opublikować inne dane. Stosujemy wówczas odpowiednio sparowane
operacje \code{release} i \code{acquire}.

\begin{lstlisting}[caption={Bezpieczna publikacja zwykłej zmiennej (C++11).}]
#include <atomic>
#include <iostream>
#include <thread>

int main() {
    int data = 0;
    std::atomic<bool> ready{false};

    std::thread producer([&] {
        data = 42;
        ready.store(true, std::memory_order_release);
    });
    std::thread consumer([&] {
        while (!ready.load(std::memory_order_acquire)) {
            std::this_thread::yield();
        }
        std::cout << data << '\n'; // 42
    });

    producer.join();
    consumer.join();
}
\end{lstlisting}

Gdy odczyt \code{acquire} zaobserwuje właściwy zapis
\code{release}, operacje \emph{synchronizes-with} gwarantują
widoczność wcześniejszego zapisu \code{data = 42}.
Nie wolno tu zamieniać obu porządków na \code{relaxed}.

\begin{center}
\begin{tabularx}{\textwidth}{@{}lX@{}}
\toprule
\textbf{Porządek} & \textbf{Zasadnicza gwarancja} \\
\midrule
\code{relaxed} & Atomowość bez synchronizacji pozostałych danych. \\
\code{release} & Publikowanie wcześniejszych efektów do odpowiedniego odczytu. \\
\code{acquire} & Odczyt odbierający efekty opublikowane przez \code{release}. \\
\code{acq_rel} & Połączenie obu kierunków dla operacji odczyt--modyfikacja--zapis. \\
\code{seq_cst} & Dodatkowy wspólny porządek operacji sekwencyjnie spójnych. \\
\bottomrule
\end{tabularx}
\end{center}

Istnieje również \code{memory_order_consume}, ale jego historia i obsługa
przez kompilatory są skomplikowane; nie wykorzystujemy go w przykładach.
Programiście początkującemu zalecamy domyślny \code{seq_cst}, a słabsze
porządki dopiero po uzasadnieniu poprawności.

\subsection{Czekanie na atomik bez aktywnej pętli: C++20}

Od C++20 obiekt atomowy udostępnia \code{wait(old)} oraz
\code{notify_one()/notify_all()}. Oczekiwanie kończy się, kiedy
zaobserwowana wartość różni się od \code{old}; wątek nie musi stale
wykonywać aktywnego odpytywania.

\begin{lstlisting}[caption={Atomowe czekanie i wybudzanie (C++20).}]
#include <atomic>
#include <iostream>
#include <thread>

int main() {
    std::atomic<bool> ready{false};
    int payload = 0;

    std::jthread worker([&] {
        payload = 42;
        ready.store(true, std::memory_order_release);
        ready.notify_one();
    });

    ready.wait(false, std::memory_order_acquire);
    std::cout << payload << '\n'; // 42
}
\end{lstlisting}

\nota{Atomowy znacznik nie jest zamiennikiem mutexa chroniącego wiele
powiązanych pól struktury. Algorytmy \emph{lock-free} bywają trudniejsze
od rozwiązania opartego na blokadzie; dodatkowo należy rozumieć m.in.
problem ABA i czas życia współdzielonych obiektów.}

% === Rozszerzenie: zaawansowany model pamięci C++ ===
\section{Model pamięci C++: formalne podstawy synchronizacji}
\label{sec:advanced-memory-model}

Sekcja~\ref{sec:race} wprowadziła pojęcie wyścigu danych. Teraz sprecyzujemy,
\emph{dlaczego} poprawny program wielowątkowy nie może polegać na fizycznym
porządku wykonywania instrukcji. Model pamięci C++ określa dopuszczalne
obserwacje programu również wtedy, gdy kompilator zmienia kolejność instrukcji,
procesor używa buforów zapisów, a rdzenie mają prywatne pamięci podręczne.
Nie jest to model jednej konkretnej architektury (np. x86 albo ARM).

\subsection{Trzy relacje: \emph{sequenced-before}, \emph{synchronizes-with} i \emph{happens-before}}

Niech $A$ i $B$ oznaczają dwa zdarzenia (ewaluacje) programu.
\begin{description}[style=nextline]
\item[\emph{Sequenced-before}]
Relacja porządkująca ewaluacje wewnątrz jednego wątku, wynikająca z zasad
sekwencjonowania w C++. Nie wolno jej utożsamiać z samą kolejnością
występowania tokenów w źródle: nie każda para podwyrażeń jest sekwencjonowana.
\item[\emph{Synchronizes-with}]
Relacja między odpowiednimi zdarzeniami synchronizacyjnymi, np. skutecznym
odblokowaniem mutexa i późniejszym jego zablokowaniem albo zapisem
\code{release} i odczytem \code{acquire}, który obserwuje zapis
(lub odpowiedni element jego sekwencji \emph{release}).
\item[\emph{Happens-before}]
Relacja logicznego porządku zdarzeń. W przykładach bez osobliwego
porządku \emph{consume} można ją konstruować z dwóch poprzednich
relacji i przechodniości. Jeżeli $A$ \emph{happens-before} $B$,
skutki $A$ ograniczają dopuszczalne obserwacje zdarzenia $B$.
\end{description}

Dla trzech zdarzeń reguła przechodniości ma postać
\[
 A\mathrel{\mathrm{hb}}B\ \land\ B\mathrel{\mathrm{hb}}C
 \quad\Longrightarrow\quad A\mathrel{\mathrm{hb}}C.
\]

Przykładowy łańcuch synchronizacji można zapisać jako
\[
\begin{aligned}
A:\text{ zapis }x
  &\xrightarrow{\mathrm{sequenced\mbox{-}before}}
    B:\text{ store(release)},\\
B&\xrightarrow{\mathrm{synchronizes\mbox{-}with}}
    C:\text{ load(acquire)},\\
C&\xrightarrow{\mathrm{sequenced\mbox{-}before}}
    D:\text{ odczyt }x.
\end{aligned}
\]
Z przechodniości wynika $A\mathrel{\mathrm{hb}}D$. Istotne jest jednak,
że operacja \code{load(acquire)} musi rzeczywiście odczytać wartość pochodzącą
z odpowiedniego zapisu \code{release} (lub jego sekwencji); samo użycie
obu etykiet w dwóch wątkach nie wystarcza.

\subsection{Wyścig danych a wyścig logiczny: przykład}

\begin{lstlisting}[caption={Dwa atomiki eliminują data race, ale nie ustalają zwycięzcy.}]
#include <atomic>
#include <thread>

int main() {
    std::atomic<int> winner{0};
    std::jthread a([&] {
        int expected = 0;
        winner.compare_exchange_strong(expected, 1);
    });
    std::jthread b([&] {
        int expected = 0;
        winner.compare_exchange_strong(expected, 2);
    });
    a.join();
    b.join();
    // winner == 1 albo winner == 2.
}
\end{lstlisting}

Wynik zależy od harmonogramu, ale program nie ma wyścigu danych.
\emph{Race condition} jest błędem dopiero wtedy, gdy program wymaga
konkretnego zwycięzcy albo narusza swój kontrakt w części harmonogramów.

\subsection{Porządek modyfikacji pojedynczego atomika}

Każdy obiekt atomowy ma własny, spójny dla wszystkich wątków
\emph{modification order}: całkowity porządek jego atomowych modyfikacji.
Nie oznacza to wspólnego globalnego porządku wszystkich atomików.
Dla modyfikacji \code{x} oraz \code{y} mogą więc istnieć obserwacje,
które intuicyjnie zaskakują, mimo że program nie zawiera \emph{data race}.

W poniższym przykładzie obie wartości startują od zera:
\begin{lstlisting}[caption={Litmus test: dwa zapisy, dwa odczyty (fragment dwóch wątków).}]
std::atomic<int> x{0}, y{0};
int r1 = -1, r2 = -1;

// Watek A:
x.store(1, std::memory_order_relaxed);
r1 = y.load(std::memory_order_relaxed);

// Watek B:
y.store(1, std::memory_order_relaxed);
r2 = x.load(std::memory_order_relaxed);
\end{lstlisting}

Przy porządku \code{relaxed} \emph{dopuszczalne} jest $r_1=r_2=0$.
Zastąpienie \emph{wszystkich czterech} operacji przez \code{seq_cst}
wyklucza ten wynik: jeden wspólny porządek musiałby wtedy spełniać
cykliczne sprzeczne ograniczenia. Test uruchomiony tylko na jednym
procesorze niczego tu nie dowodzi: brak obserwacji nie oznacza zakazu
w standardzie.

\subsection{Porządki atomowe w praktyce}

\begin{center}
\small
\begin{tabularx}{\textwidth}{@{}p{0.25\textwidth}X@{}}
\toprule
\textbf{Porządek} & \textbf{Praktyczny sens i typowe zastosowanie} \\
\midrule
\code{relaxed} & Niepodzielność i porządek modyfikacji danego atomika; liczniki statystyk, gdy nie synchronizują innego stanu. \\
\code{release} & Publikuje efekty poprzedzających operacji; stosowany do zapisów lub odpowiednich operacji RMW. \\
\code{acquire} & Odbiera publikację po odczytaniu właściwej wartości; stosowany do odczytów lub odpowiednich RMW. \\
\code{acq_rel} & Oba kierunki synchronizacji podczas operacji typu read--modify--write. \\
\code{seq_cst} & Domyślny, mocniejszy porządek; dodatkowo wspólne uporządkowanie wszystkich operacji \code{seq_cst}. \\
\bottomrule
\end{tabularx}
\end{center}

\code{memory_order_consume} jest historycznie problematyczny; nie stanowi
dobrego punktu wyjścia do pisania przenośnego kodu. Wybór słabszego
porządku należy poprzedzić argumentem poprawności, a nie samym wynikiem
benchmarku. Nie wszystkie porządki są dozwolone dla każdej operacji,
np. zwykły \code{store} nie może użyć \code{memory_order_acquire}.

\subsection{Publikacja danych i sekwencja \emph{release}}

Rozpatrzmy trzy zdarzenia: producent wypełnia dane, drugi wątek
zwiększa atomowy licznik, a trzeci odbiera wynik. W C++20 (i wcześniej)
sekwencja \emph{release} po zapisie \code{release} może być kontynuowana
przez kolejne atomowe operacje \emph{read--modify--write} tego samego
atomika, także wykonywane przez inne wątki. Dlatego odczyt \code{acquire}
wartości zapisanej przez taką operację RMW może synchronizować się
z początkowym zapisem \code{release}.

\begin{lstlisting}[caption={Sekwencja release: fragment przy ustalonej kolejności operacji.}]
int payload = 0;
std::atomic<int> state{0};

// Producent (A):
payload = 42;
state.store(1, std::memory_order_release);

// Posrednik (B), DOPIERO po wykonaniu powyzszego store:
state.fetch_add(1, std::memory_order_relaxed); // 1 -> 2

// Konsument (C): odczytuje 2 z RMW posrednika:
if (state.load(std::memory_order_acquire) == 2) {
    // Wolno odczytac payload == 42.
}
\end{lstlisting}

Warunek o kolejności jest częścią przykładu: bez niego pośrednik mógłby
zwiększyć \code{state} jeszcze przed publikacją. Zwykły zapis wykonany
przez inny wątek przerywa tę sekwencję; nie należy jej mylić z dowolną
serią zapisów do tego samego atomika.

\subsection{Operacje CAS: sukces i porażka}

\code{compare_exchange_weak} i \code{compare_exchange_strong} mają
przeciążenia z osobnymi porządkami pamięci dla sukcesu i porażki.
Operacja nieudana jest \emph{odczytem}, a nie zapisem. Dlatego jej porządek
nie może być \code{release} ani \code{acq_rel} oraz nie może być silniejszy
od porządku sukcesu w sensie wymagań standardu.

\begin{lstlisting}[caption={Pętla CAS: zwiększanie licznika, przykład praktyczny.}]
#include <atomic>

void add_one(std::atomic<int>& x) {
    int expected = x.load(std::memory_order_relaxed);
    while (!x.compare_exchange_weak(
        expected, expected + 1,
        std::memory_order_acq_rel,
        std::memory_order_acquire)) {
        // expected zostalo zaktualizowane po niepowodzeniu.
    }
}
\end{lstlisting}

W tym przypadku \code{fetch_add(1)} jest prostsze i zwykle lepsze.
CAS ma sens zwłaszcza przy warunkowej zmianie stanu lub złożonych
strukturach atomowych. \code{weak} może nie powieść się pozornie,
więc jest wygodny w pętli; \code{strong} lepiej nadaje się do
pojedynczych prób. Użycie \code{acq_rel} w powyższym liczniku służy
zilustrowaniu API; gdy licznik niczego nie publikuje, wystarczy
\code{memory_order_relaxed}.

\subsection{Bariery pamięci: \texttt{atomic\_thread\_fence}}

\code{std::atomic_thread_fence(order)} nakłada ograniczenia porządkowania
pamięci, ale \emph{sama} bariera nie synchronizuje magicznie dwóch wątków.
Potrzebna jest odpowiednia operacja atomowa, która połączy strony
synchronizacji. Przykład celowo przenosi \code{release}/\code{acquire}
z atomika na bariery:

\begin{lstlisting}[caption={Publikacja przez dwie bariery i atomowy znacznik (C++11).}]
#include <atomic>
#include <thread>
#include <cassert>

int main() {
    int payload = 0;
    std::atomic<bool> ready{false};

    std::thread producer([&] {
        payload = 42;
        std::atomic_thread_fence(std::memory_order_release);
        ready.store(true, std::memory_order_relaxed);
    });
    std::thread consumer([&] {
        while (!ready.load(std::memory_order_relaxed)) {
            std::this_thread::yield();
        }
        std::atomic_thread_fence(std::memory_order_acquire);
        assert(payload == 42);
    });
    producer.join();
    consumer.join();
}
\end{lstlisting}

Dla tego układu bariery łączą się poprzez odczyt \code{ready}, który
obserwuje zapis producenta. Zwykłe \code{ready.store(release)} oraz
\code{ready.load(acquire)} byłoby prostsze. \code{atomic_signal_fence}
ma inny cel: ogranicza określone przestawienia przez kompilator,
a nie pełni roli zwykłej synchronizacji między wątkami.

\subsection{Atomowość, blokady i gwarancje postępu}

\begin{description}[style=nextline]
\item[\emph{Blocking}] Wątek może oczekiwać na zwolnienie blokady przez
innego uczestnika, który np. został wywłaszczony przez system.
\item[\emph{Lock-free}] W standardowym ujęciu algorytmowym system jako
całość może robić postęp mimo tego, że pewne wątki wielokrotnie ponawiają
operację. Pojedynczy wątek może ulegać zagłodzeniu.
\item[\emph{Wait-free}] Każdy uczestnik kończy operację w ograniczonej
liczbie własnych kroków, niezależnie od harmonogramu innych uczestników.
\end{description}

\code{atomic<T>::is_lock_free()} bada własność implementacji operacji
na tym atomiku; nie jest dowodem, że dowolny algorytm zbudowany z
atomików jest \emph{lock-free}. Nie wolno też wnioskować, że
\emph{lock-free} zawsze oznacza szybszy kod. Rywalizacja o jedną
linię cache może dominować koszt obliczeń.

\subsection{Problem ABA i odzyskiwanie pamięci}

Niech wierzchołek współbieżnego stosu zawiera wskaźnik \code{A}.
Wątek~1 odczytuje \code{A} i zapamiętuje wskaźnik do następnego
węzła, po czym zostaje wstrzymany. Wątek~2 usuwa \code{A}, usuwa
kolejny element \code{B} i ponownie umieszcza \code{A} na szczycie.
Wątek~1 widzi znowu \code{A}: porównanie wskaźnika przez CAS może się
powieść, choć pomiędzy odczytem i zapisem struktura zmieniła się
istotnie. Co gorsza, wcześniej zapamiętany węzeł mógł zostać już
zwolniony. To właśnie \emph{problem ABA}.

\[
\boxed{\text{A}\to\text{B}\to\text{C}}
\quad\longrightarrow\quad
\boxed{\text{B}\to\text{C}}
\quad\longrightarrow\quad
\boxed{\text{A}\to\text{C}}.
\]

Do typowych narzędzi należą: oznaczanie wskaźnika licznikiem wersji
(\emph{tagged pointer}), \emph{hazard pointers}, \emph{epoch-based
reclamation} i \emph{reference counting}. Zabezpieczenie samego
porównania wskaźników nie załatwia problemu \emph{czasu życia obiektu}.
Dla zwykłej kolejki aplikacyjnej \code{std::mutex} bywa
najprostszy i najbardziej wiarygodny.

\subsection{Zagłodzenie, \emph{livelock} i błędne podwójne sprawdzenie}

\emph{Deadlock} oznacza brak postępu wskutek cyklicznego oczekiwania.
\emph{Livelock} występuje, gdy wątki aktywnie reagują na siebie,
ale nie kończą pracy. \emph{Starvation} to brak postępu konkretnego
wątku mimo że inne wątki pracują. Te trzy zjawiska należy odróżniać.

W starszym kodzie spotyka się tzw. \emph{double-checked locking}.
Bez poprawnych atomików albo mutexa nie można go naprawić przez samo
\code{volatile}. W nowym kodzie znacznie lepiej używać bezpiecznej
inicjalizacji statycznej zmiennej lokalnej (C++11) lub \code{std::call_once}.
\code{volatile} nie zapewnia atomowości ani synchronizacji wątków.



\section{Wynik przyszłego obliczenia: \texttt{future} i \texttt{promise}}
\label{sec:future-promise}

\code{std::future<T>} reprezentuje możliwość odebrania \emph{jednego}
wyniku typu \code{T}, który może jeszcze nie być dostępny.
\code{std::promise<T>} umożliwia producentowi ten wynik ustawić.
Te obiekty odnoszą się do wspólnego \emph{stanu asynchronicznego}
(ang. \emph{shared state}). Sam \code{promise} nie uruchamia wątku.

\begin{center}
\fbox{\small\strut\code{promise.set_value(x)}\quad
$\longrightarrow$\quad stan współdzielony\quad $\longrightarrow$\quad
\code{future.get()}}
\end{center}

\begin{lstlisting}[caption={Ręczne przekazanie wyniku z wątku (C++11).}]
#include <future>
#include <iostream>
#include <thread>
#include <utility>

void compute(std::promise<int> producer) {
    producer.set_value(6 * 7);
}

int main() {
    std::promise<int> promise;
    std::future<int> result = promise.get_future();

    std::thread worker(compute, std::move(promise));
    std::cout << result.get() << '\n'; // 42
    worker.join();
}
\end{lstlisting}

\code{future.get()} czeka, jeśli wynik nie jest jeszcze gotowy, po czym
odbiera go (lub ponownie zgłasza wyjątek). Zwykłe \code{future}
jest jednokrotnego odbioru: po \code{get()} obiekt przestaje być
ważnym uchwytem do stanu i ponowne \code{get()} nie jest dozwolone.
\code{wait()} tylko czeka, natomiast \code{wait_for()} i \code{wait_until()}
pozwalają ograniczyć czas oczekiwania.

\subsection{Przekazywanie wyjątku i \emph{broken promise}}

Producent może przekazać wyjątek za pomocą \code{set_exception()}.
Jeśli \code{promise} zostanie zniszczony bez przekazania wyniku,
stan otrzymuje błąd \code{broken_promise}, który zostanie zgłoszony
podczas odbioru przez \code{future.get()}.

\begin{lstlisting}[caption={Przekazanie wyjątku przez promise (C++14).}]
#include <exception>
#include <future>
#include <iostream>
#include <stdexcept>
#include <thread>
#include <utility>

int main() {
    std::promise<int> producer;
    auto f = producer.get_future();

    std::thread worker([p = std::move(producer)]() mutable {
        try {
            throw std::runtime_error("Calculation failed");
        } catch (...) {
            p.set_exception(std::current_exception());
        }
    });

    try {
        std::cout << f.get() << '\n';
    } catch (const std::exception& e) {
        std::cout << e.what() << '\n';
    }
    worker.join();
}
\end{lstlisting}

\subsection{Wiele odbiorców: \texttt{shared\_future}}

\code{std::shared_future<T>} można kopiować i wielokrotnie odczytywać.
Każdy odbiorca uzyskuje dostęp do tego samego wyniku. Uzyskujemy go
np. przez \code{future.share()}.

\begin{lstlisting}[caption={Dwa wątki odczytują ten sam wynik (C++11).}]
#include <future>
#include <iostream>
#include <thread>

int main() {
    std::promise<int> p;
    std::shared_future<int> f = p.get_future().share();
    int a = 0, b = 0;

    std::thread first([&] { a = f.get(); });
    std::thread second([&] { b = f.get(); });

    p.set_value(42);
    first.join();
    second.join();
    std::cout << a << ' ' << b << '\n'; // 42 42
}
\end{lstlisting}

\section{Uruchamianie zadań: \texttt{std::async}}
\label{sec:async}

\code{std::async} wiąże wywołanie funkcji ze stanem asynchronicznym
i zwraca \code{std::future<R>} dla typu wyniku \code{R}.
W odróżnieniu od \code{promise} nie trzeba samodzielnie obsługiwać
przekazywania wartości ani wyjątków. Wariant
\code{std::launch::async} uruchamia funkcję w innym wątku.

\begin{lstlisting}[caption={Dwa niezależne obliczenia przez async (C++11).}]
#include <future>
#include <iostream>

int square(int x) { return x * x; }

int main() {
    auto f1 = std::async(std::launch::async, square, 6);
    auto f2 = std::async(std::launch::async, square, 7);

    // Oba zadania wystartowaly przed pobraniem wynikow.
    std::cout << f1.get() + f2.get() << '\n'; // 85
}
\end{lstlisting}

\subsection{Polityki uruchamiania}

\begin{center}
\begin{tabularx}{\textwidth}{@{}lX@{}}
\toprule
\textbf{Polityka} & \textbf{Znaczenie} \\
\midrule
\code{std::launch::async} & Funkcja zostaje uruchomiona w osobnym wątku; wywołujący może kontynuować pracę. \\
\code{std::launch::deferred} & Funkcja zostaje wykonana dopiero przy pierwszym nieczasowym oczekiwaniu (np. \code{get()} lub \code{wait()}), w wątku, który tę operację wykonuje. \\
Brak jawnej polityki & Implementacja może wybrać \code{async}, \code{deferred} (lub dopuszczoną politykę implementacyjną). \\
\bottomrule
\end{tabularx}
\end{center}

Nie wolno utożsamiać samego wywołania \code{std::async(f)}
z gwarancją równoległości. Jeśli potrzebujesz niezależnego wątku,
wybierz jawnie \code{std::launch::async}.

\begin{lstlisting}[caption={Opóźnione wykonanie: deferred (C++11).}]
#include <future>
#include <iostream>

int main() {
    auto f = std::async(std::launch::deferred, [] {
        std::cout << "Executing the task\n";
        return 42;
    });
    std::cout << "Before get\n";
    std::cout << f.get() << '\n';
}
\end{lstlisting}

W tym przypadku komunikat \code{Before get} zawsze poprzedza wykonanie
funkcji odroczonej, ponieważ została ona uruchomiona dopiero przez \code{get()}.

\subsection{Destruktor tymczasowego \texttt{future} może blokować}

Istotna pułapka polega na tym, że zniszczenie \emph{ostatniego} uchwytu
do niezakończonego stanu utworzonego przez \code{std::async}
z polityką \code{async} może czekać na ukończenie zadania.
Zapis poniżej nie uruchamia efektywnie obu prac obok siebie:

\begin{lstlisting}[caption={Pułapka: tymczasowe future; fragment ilustracyjny.}]
// Kazdy tymczasowy future jest niszczony na koncu instrukcji.
std::async(std::launch::async, expensive_a);
std::async(std::launch::async, expensive_b);
\end{lstlisting}

Poprawnie przechowujemy wyniki obu wywołań, zanim zaczniemy czekać:

\begin{lstlisting}[caption={Zadania mogą pracować równolegle; fragment.}]
auto a = std::async(std::launch::async, expensive_a);
auto b = std::async(std::launch::async, expensive_b);
a.wait();
b.wait();
\end{lstlisting}

\subsection{Oczekiwanie z limitem czasu i polityka deferred}

\code{wait_for()} zwraca jeden z trzech statusów:
\code{ready}, \code{timeout} lub \code{deferred}.
Przy zadaniu odroczonym samo \code{wait_for()} nie uruchamia pracy.
Niepoprawna pętla sprawdzająca wyłącznie \code{ready}
może więc kręcić się bez końca.

\begin{lstlisting}[caption={Rozróżnianie statusów future (C++11).}]
#include <chrono>
#include <future>
#include <iostream>

int main() {
    auto f = std::async(std::launch::deferred, [] { return 42; });
    auto status = f.wait_for(std::chrono::milliseconds(1));

    if (status == std::future_status::deferred) {
        std::cout << "Deferred: get starts the computation\n";
    } else if (status == std::future_status::timeout) {
        std::cout << "Not ready yet\n";
    } else {
        std::cout << "Already ready\n";
    }
    std::cout << f.get() << '\n';
}
\end{lstlisting}

\subsection{Wyjątki z zadań i brak automatycznego anulowania}

Jeśli funkcja uruchomiona przez \code{std::async} zgłosi wyjątek,
wyjątek jest przechowywany w stanie asynchronicznym, a \code{get()}
zgłasza go w wątku odbierającym wynik. \code{wait_for()}
\emph{nie anuluje} pracy przy przekroczeniu czasu: tylko informuje,
że wynik jeszcze nie jest gotowy.

\begin{lstlisting}[caption={Wyjątek z async zostaje przekazany przez get() (C++11).}]
#include <future>
#include <iostream>
#include <stdexcept>

int main() {
    auto f = std::async(std::launch::async, []() -> int {
        throw std::runtime_error("Task failed");
    });
    try {
        std::cout << f.get() << '\n';
    } catch (const std::runtime_error& e) {
        std::cout << e.what() << '\n';
    }
}
\end{lstlisting}

\section{\texttt{std::packaged\_task}: funkcja z wynikiem future}
\label{sec:packaged}

\code{std::packaged_task<R(Args...)>} opakowuje funkcję w zadanie,
którego wynik lub wyjątek trafia do powiązanego \code{future}.
W odróżnieniu od \code{std::async} samo utworzenie
\code{packaged_task} nie uruchamia żadnego wątku.

\begin{lstlisting}[caption={Packaged task i ręcznie uruchomiony wątek (C++11).}]
#include <future>
#include <iostream>
#include <thread>
#include <utility>

int main() {
    std::packaged_task<int(int)> task([](int x) {
        return x * x;
    });
    auto f = task.get_future();
    std::thread worker(std::move(task), 7);

    std::cout << f.get() << '\n'; // 49
    worker.join();
}
\end{lstlisting}

To jeden z elementów potrzebnych do zbudowania puli wątków:
\emph{task} trafia do kolejki, wątek roboczy go pobiera,
a klient otrzymuje \code{future}. Pula wątków jest szczególnie
przydatna dla bardzo wielu małych zadań, ponieważ tworzenie
odrębnego wątku dla każdego zadania może być kosztowne.

\section{\texttt{std::jthread} i kooperacyjne zatrzymywanie (C++20)}
\label{sec:jthread}

\code{std::jthread} jest obiektem zarządzającym wątkiem zgodnie z RAII.
Jeżeli jest dołączalny, destruktor żąda zatrzymania przez
\code{request_stop()}, a następnie wywołuje \code{join()}.
Żądanie zatrzymania nie wymusza przerwania funkcji: wątek musi
sam sprawdzać odpowiedni \code{stop_token}.

\begin{lstlisting}[caption={Kooperacyjne zatrzymywanie wątku (C++20).}]
#include <chrono>
#include <iostream>
#include <thread>

int main() {
    std::jthread worker([](std::stop_token stop) {
        while (!stop.stop_requested()) {
            // Przykladowa porcja pracy.
            std::this_thread::sleep_for(std::chrono::milliseconds(5));
        }
        std::cout << "Stopped\n";
    });

    std::this_thread::sleep_for(std::chrono::milliseconds(25));
    worker.request_stop();
    // Destruktor worker dolaczy watek.
}
\end{lstlisting}

Można również przekazywać \code{stop_token} do zadań, które
reagują na anulowanie w większej liczbie punktów. Dla złożonych
kolejek C++20 oferuje \code{condition_variable_any} z przeciążeniem
obsługującym żądanie zatrzymania; zwykłe \code{condition_variable}
nie ma analogicznego przeciążenia ze \code{stop_token}.

\section{Inne przydatne mechanizmy biblioteki standardowej}

\subsection{\texttt{shared\_mutex}: wielu czytelników, jeden pisarz}

\code{std::shared_mutex} (C++17) pozwala wielu czytelnikom jednocześnie
uzyskać blokadę współdzieloną, a pisarzowi --- blokadę wyłączną.
Czytanie chronimy przez \code{std::shared_lock}, zapisy przez
\code{std::unique_lock}. Nie gwarantuje to braku zagłodzenia wątków;
polityka przydziału blokad może zależeć od implementacji.

\begin{lstlisting}[caption={Blokada współdzielona i wyłączna (C++17).}]
#include <iostream>
#include <mutex>
#include <shared_mutex>

int main() {
    std::shared_mutex m;
    int value = 10;
    {
        std::shared_lock<std::shared_mutex> read_lock(m);
        std::cout << value << '\n';
    }
    {
        std::unique_lock<std::shared_mutex> write_lock(m);
        value = 20;
    }
}
\end{lstlisting}

\subsection{\texttt{call\_once} i inicjalizacja jednokrotna}

\code{std::call_once} oraz \code{std::once_flag} pozwalają
wykonać funkcję inicjalizującą co najwyżej raz pomyślnie.
Jeżeli inicjalizacja rzuci wyjątek, późniejsze wywołanie może
ponowić próbę. Lokalna statyczna zmienna ma w C++11
bezpieczną współbieżnie inicjalizację, co często jest prostszą alternatywą.

\begin{lstlisting}[caption={Inicjalizacja wykonywana raz (C++11).}]
#include <iostream>
#include <mutex>
#include <thread>

std::once_flag init_flag;
int shared_value = 0;

void initialize() { shared_value = 42; }
void work() {
    std::call_once(init_flag, initialize);
    // Kazdy watek po call_once widzi zakonczona inicjalizacje.
}

int main() {
    std::thread a(work), b(work);
    a.join();
    b.join();
    std::cout << shared_value << '\n';
}
\end{lstlisting}

\subsection{\texttt{thread\_local}: osobny egzemplarz dla wątku}

Zmienne zadeklarowane \code{thread_local} mają osobny egzemplarz
w każdym wątku. Nie rozwiązują automatycznie problemu komunikacji
między wątkami, ale przydają się np. do lokalnych buforów roboczych.

\begin{lstlisting}[caption={Pamięć lokalna dla każdego wątku (C++11).}]
#include <iostream>
#include <thread>

thread_local int local_counter = 0;
void count() { ++local_counter; }

int main() {
    std::thread a(count), b(count);
    a.join();
    b.join();
    std::cout << local_counter << '\n'; // 0 w watku main
}
\end{lstlisting}

\subsection{Semafory, zapadki i bariery w C++20}

\begin{description}[style=nextline]
\item[\code{std::counting_semaphore<N>}]
Ogranicza liczbę równoczesnych wejść. Operacja \code{acquire()}
pobiera zezwolenie, a \code{release()} je zwraca.
\item[\code{std::latch}]
Jednorazowy punkt synchronizacji: po odpowiedniej liczbie zgłoszeń
\code{count_down()} czekające wątki mogą kontynuować.
\item[\code{std::barrier}]
Punkt synchronizacji wielokrotnego użytku dla kolejnych faz pracy.
\end{description}

\begin{lstlisting}[caption={Semafor dopuszcza maksymalnie dwa zadania naraz (C++20).}]
#include <chrono>
#include <iostream>
#include <semaphore>
#include <thread>
#include <vector>

int main() {
    std::counting_semaphore<2> slots(2);
    std::vector<std::jthread> workers;
    for (int i = 0; i < 4; ++i) {
        workers.emplace_back([&, i] {
            slots.acquire();
            // Jednoczesnie w sekcji moze byc najwyzej 2 watki.
            std::this_thread::sleep_for(std::chrono::milliseconds(10));
            slots.release();
        });
    }
    // Destruktory jthread dolaczaja watki przed destruktorem slots.
}
\end{lstlisting}

\nota{W przykładzie z semaforem zmienne zostały zadeklarowane w kolejności,
dzięki której wektor \code{workers} zostaje zniszczony przed \code{slots}.
Ma to znaczenie dla czasu życia obiektów współdzielonych przez wątki.}

% === Rozszerzenie: zadania, pula wątków, anulowanie ===
\section{Zadania, pule wątków i planowanie pracy}
\label{sec:thread-pool}

Wątek jest zasobem systemowym, a zadanie jest jednostką pracy. Wiele
krótkich zadań nie powinno wymagać utworzenia osobnego wątku na każde
wywołanie. \emph{Thread pool} utrzymuje ograniczoną liczbę pracowników,
którzy pobierają zadania ze współdzielonej kolejki.

\subsection{Minimalna pula wątków ze zwracaniem \texttt{future}}

Poniższa implementacja jest przykładem dydaktycznym w C++20.
Łączy kilka poznanych mechanizmów:
\begin{itemize}
\item \code{std::jthread} i \code{std::condition_variable} do obsługi pracowników;
\item \code{std::packaged_task} i \code{std::future} do przekazywania wyników;
\item \code{std::function} do przechowywania zadań w jednej kolejce.
\end{itemize}
Metoda \code{submit} zwraca \code{future} z wynikiem.
Wyjątek zgłoszony przez zadanie zostaje ponowiony podczas \code{get()}.

\begin{lstlisting}[caption={Prosta pula: kompletna implementacja C++20.}]
#include <condition_variable>
#include <cstddef>
#include <functional>
#include <future>
#include <memory>
#include <mutex>
#include <queue>
#include <stdexcept>
#include <thread>
#include <type_traits>
#include <utility>
#include <vector>

class ThreadPool {
    std::mutex m_;
    std::condition_variable cv_;
    std::queue<std::function<void()>> tasks_;
    bool closing_ = false;
    // workers_ jest ostatni: jego destruktory wywolaja join
    // PRZED zniszczeniem m_, cv_ i kolejki.
    std::vector<std::jthread> workers_;

public:
    explicit ThreadPool(std::size_t count) {
        if (count == 0) throw std::invalid_argument("count=0");
        workers_.reserve(count);
        for (std::size_t i = 0; i < count; ++i) {
            workers_.emplace_back([this] {
                for (;;) {
                    std::function<void()> work;
                    {
                        std::unique_lock lock(m_);
                        cv_.wait(lock, [&] {
                            return closing_ || !tasks_.empty();
                        });
                        if (closing_ && tasks_.empty()) return;
                        work = std::move(tasks_.front());
                        tasks_.pop();
                    }
                    work(); // poza sekcja krytyczna
                }
            });
        }
    }

    ThreadPool(const ThreadPool&) = delete;
    ThreadPool& operator=(const ThreadPool&) = delete;

    template<class F>
    auto submit(F&& f)
        -> std::future<std::invoke_result_t<std::decay_t<F>&>> {
        using Result = std::invoke_result_t<std::decay_t<F>&>;
        auto task = std::make_shared<std::packaged_task<Result()>>(
            std::forward<F>(f));
        auto result = task->get_future();
        {
            std::lock_guard lock(m_);
            if (closing_)
                throw std::runtime_error("Pool is closing");
            tasks_.emplace([task] { (*task)(); });
        }
        cv_.notify_one();
        return result;
    }

    ~ThreadPool() {
        {
            std::lock_guard lock(m_);
            closing_ = true;
        }
        cv_.notify_all();
        // Destruktory jthread czekaja na pracownikow.
    }
};
\end{lstlisting}

\begin{lstlisting}[caption={Użycie puli i odbiór wyników; dołącz do programu z klasą ThreadPool.}]
#include <iostream>

int main() {
    ThreadPool pool(3);
    auto a = pool.submit([] { return 40 + 2; });
    auto b = pool.submit([] { return 6 * 7; });
    std::cout << a.get() << " " << b.get() << '\n';
}
\end{lstlisting}

Przykład nie implementuje wszystkich funkcji puli produkcyjnej:
nie ma limitu długości kolejki, priorytetów, \emph{work stealing},
możliwości natychmiastowego porzucenia zadania ani kooperacyjnego
anulowania pojedynczego zadania. Destruktor wykonuje model
\emph{graceful shutdown}: przestaje przyjmować zadania, pozwala
dokończyć prace oczekujące w kolejce i czeka na zakończenie pracowników.
Wywołujący musi zadbać, żeby inne wątki nie używały obiektu puli podczas
jego niszczenia. Nie należy też niszczyć puli z jej własnego pracownika.

\nota{Gdy pula ma jednego pracownika, zadanie oczekujące synchronicznie
na \code{future} innego zadania umieszczonego w \emph{tej samej} puli
może wywołać zakleszczenie. Równie ważna jest kolejność niszczenia
składowych wykorzystywanych przez pracowników.}

\subsection{Anulowanie zadania: prośba, nie wymuszone przerwanie}

\code{std::stop_source} dostarcza możliwość wysłania żądania, a
\code{std::stop_token} pozwala je obserwować. Mogą istnieć liczni
odbiorcy jednego tokenu, a sygnał pozostaje ustawiony po wywołaniu
\code{request_stop()}. Mechanizm sam nie przerywa instrukcji,
blokującego wywołania systemowego ani obliczenia, które ignoruje token.

\begin{lstlisting}[caption={Jedna prośba o zatrzymanie i kilka punktów kontroli.}]
#include <chrono>
#include <iostream>
#include <stop_token>
#include <thread>

void work(std::stop_token token) {
    for (int i = 0; i < 100; ++i) {
        if (token.stop_requested()) {
            std::cout << "Cancelled after " << i << " steps\n";
            return;
        }
        std::this_thread::sleep_for(std::chrono::milliseconds(1));
    }
}

int main() {
    std::stop_source source;
    std::jthread worker(work, source.get_token());
    std::this_thread::sleep_for(std::chrono::milliseconds(5));
    source.request_stop();
}
\end{lstlisting}

Możliwa jest także reakcja w callbacku rejestrowanym przez
\code{std::stop_callback}. Trzeba uważać, by callback nie blokował
mutexa, który jest już trzymany przez wątek wywołujący
\code{request_stop()}, ponieważ callback może zostać wykonany
\emph{synchronicznie} w tym wywołaniu.

\subsection{Czekanie z możliwością zatrzymania}

\code{std::condition_variable_any} ma w C++20 przeciążenie
\code{wait(lock, token, pred)}. Dzięki niemu kolejka zadań może
wybudzać się również po prośbie o zatrzymanie.

\begin{lstlisting}[caption={Oczekiwanie na stan lub anulowanie w C++20.}]
#include <condition_variable>
#include <mutex>
#include <stop_token>

std::mutex m;
std::condition_variable_any cv;
bool ready = false;

void wait_for_work(std::stop_token token) {
    std::unique_lock lock(m);
    bool completed = cv.wait(lock, token, [&] { return ready; });
    if (!completed) {
        // Stop requested zanim predykat stal sie prawdziwy.
        return;
    }
    // W tym miejscu trzymamy m, a ready jest prawdziwe.
}
\end{lstlisting}

Sam token nie zastępuje sygnalizacji \code{cv.notify_one()} po zmianie
predykatu. Wykorzystanie tokenów nie zwalnia także z wymagania, by
stan zadania i jego zasoby pozostawały ważne aż do zakończenia pracy.

\subsection{Wiele zadań i wyjątki: kiedy użyć \texttt{async}, a kiedy puli?}

\code{std::async(std::launch::async, f)} jest dobry, gdy liczba zadań
jest mała, a każde zadanie względnie kosztowne. Pula sprawdza się,
gdy zadań jest wiele i nie chcemy tworzyć osobnego wątku dla każdego.
Z kolei \code{std::jthread} jest wygodny, gdy potrzebujemy długowiecznego
pracownika wykonującego pętlę. W praktyce można też użyć gotowego
\emph{executora} lub biblioteki wyspecjalizowanej w planowaniu zadań.

Gdy funkcja w wątku rzuca wyjątek i nie przechwytuje go lokalnie,
program może zakończyć się przez \code{std::terminate}. Wymianę
wyjątków między wątkami organizujemy przez \code{future},
\code{promise::set_exception} lub \code{std::exception_ptr}.

\section{Skalowanie i koszt współbieżności}

Przyspieszenie równoległego algorytmu jest ograniczone nie tylko przez
liczbę rdzeni, ale także przez odsetek części sekwencyjnej i koszty
synchronizacji. Jeśli część $s\in[0,1]$ pracy musi pozostać sekwencyjna,
to idealizowane prawo Amdahla daje
\[
S(p)\le \frac{1}{s+(1-s)/p},
\]
gdzie $p$ to liczba procesorów/rdzeni roboczych. Jest to uproszczony
model pomijający np. pamięć, uruchamianie wątków i nierównomierność pracy.

W praktyce należy rozważyć: narzut utworzenia i obudzenia wątku,
przeciążenie wspólnego mutexa, współdzielenie linii cache
(\emph{false sharing}), ograniczoną przepustowość pamięci i liczbę
rzeczywiście fizycznych rdzeni. \code{std::thread::hardware_concurrency()}
może być użyteczną wskazówką, ale nie gwarantuje liczby dostępnych
rdzeni i może zwrócić zero.

Równoległość nie jest synonimem wydajności: drobne zadania,
trudne do przewidzenia skoki oraz intensywna konkurencja o pamięć
mogą działać wolniej od wariantu jednowątkowego. Pomiary należy
wykonywać na reprezentatywnych danych w konfiguracji Release,
a poprawność sprawdzać niezależnie od benchmarków.



\section{Przykład numeryczny: równoległe sumowanie}
\label{sec:numerical}

Załóżmy, że mamy duży wektor wartości i chcemy policzyć sumę. Możemy
podzielić zakres na dwie rozłączne części, jedną zsumować w zadaniu
\code{std::async}, a drugą w bieżącym wątku.

\begin{lstlisting}[caption={Dwa niezależne fragmenty obliczenia sumy (C++11).}]
#include <future>
#include <iostream>
#include <numeric>
#include <vector>

int main() {
    std::vector<double> values(100000, 1.0);
    auto middle = values.begin() + values.size() / 2;

    auto right = std::async(std::launch::async, [&] {
        return std::accumulate(middle, values.end(), 0.0);
    });

    double left = std::accumulate(values.begin(), middle, 0.0);
    double total = left + right.get();
    std::cout << total << '\n'; // 100000
}
\end{lstlisting}

Odczytywane fragmenty \code{values} pozostają niezmienne aż do zakończenia
obliczeń, więc nie ma wyścigu danych. Trzeba jednak pamiętać, że
\emph{koszt uruchomienia zadania} może przewyższać korzyść, jeżeli
obliczenia są zbyt krótkie. Przy sumach zmiennoprzecinkowych inna
kolejność dodawania może ponadto zmienić wynik z powodu zaokrągleń.

% === Rozszerzenie: OpenMP ===
\section{OpenMP: programowanie równoległe przez dyrektywy}
\label{sec:openmp}

\emph{OpenMP} (\emph{Open Multi-Processing}) to standard API do
programowania równoległego na pamięci współdzielonej, stosowany
w C, C++ i Fortranie. Nie jest pojedynczą klasą ani tradycyjną
biblioteką szablonową C++: łączy \emph{dyrektywy kompilatora},
\emph{bibliotekę funkcji wykonawczych} (\code{omp.h}) i \emph{zmienne
środowiskowe}. Posiada również konstrukcje do pracy na akceleratorach,
lecz tutaj skupimy się na wielordzeniowym CPU.

Zaletą OpenMP jest możliwość stopniowego równoleglenia istniejących pętli.
Kod źródłowy pozostaje zazwyczaj podobny do wersji sekwencyjnej.
Ceną jest konieczność rozumienia prywatności zmiennych, zależności
pomiędzy iteracjami i synchronizacji.

\nota{Specyfikacja OpenMP ma własne wersje (m.in. 5.2 i 6.0),
niezależne od C++20. W przykładach poniżej używamy głównie
klasycznych dyrektyw dostępnych w wielu kompilatorach. Samo istnienie
mechanizmu w specyfikacji nie gwarantuje, że konkretna wersja GCC,
Clang, MSVC lub Intel oneAPI obsługuje go w pełni.}

\subsection{Kompilacja i konfiguracja OpenMP}

Dla GCC i typowych konfiguracji Clang można użyć flagi
\code{-fopenmp}. W Clangu konieczne może być dodatkowo
zainstalowanie odpowiedniego środowiska wykonawczego OpenMP
(np. \code{libomp}). W MSVC wsparcie i zestaw flag zależą od
wersji oraz wybranego trybu OpenMP.

\begin{lstlisting}[language=bash,caption={Kompilacja programu OpenMP w GCC.}]
g++ -std=c++20 -O2 -Wall -Wextra -fopenmp example.cpp -o example
./example

# Przykladowe uruchomienie z 4 watkami:
OMP_NUM_THREADS=4 ./example
\end{lstlisting}

W CMake można uniknąć ręcznego wpisywania flag:
\begin{lstlisting}[language=bash,caption={Minimalny CMakeLists.txt (tekst pliku).}]
cmake_minimum_required(VERSION 3.16)
project(OpenMPDemo LANGUAGES CXX)
set(CMAKE_CXX_STANDARD 20)
find_package(OpenMP REQUIRED COMPONENTS CXX)
add_executable(example example.cpp)
target_link_libraries(example PRIVATE OpenMP::OpenMP_CXX)
\end{lstlisting}

Biblioteka wykonawcza udostępnia funkcje do odczytu stanu środowiska:
\begin{itemize}
\item \code{omp_get_thread_num()} --- identyfikator wątku;
\item \code{omp_get_num_threads()} --- liczba uczestników aktualnego zespołu;
\item \code{omp_get_max_threads()} --- górna granica następnego zespołu;
\item \code{omp_set_num_threads()} --- ustawianie żądanej liczby wątków;
\item \code{omp_get_wtime()} --- czas ścienny w sekundach.
\end{itemize}
Poza regionem równoległym \code{omp_get_num_threads()} zwraca~1.
Wartość \code{omp_get_max_threads()} nie jest liczbą fizycznych rdzeni.

\subsection{Pierwszy region: \texttt{parallel}}

\code{\#pragma omp parallel} tworzy zespół wątków, którego uczestnicy
wykonują blok. Wątek napotykający dyrektywę staje się wątkiem
głównym tego zespołu (numer~0). Zakończenie bloku obejmuje domyślną
synchronizację uczestników.

\begin{lstlisting}[caption={Region równoległy OpenMP.}]
#include <iostream>
#include <omp.h>

int main() {
    #pragma omp parallel num_threads(4)
    {
        int id = omp_get_thread_num();
        int count = omp_get_num_threads();
        #pragma omp critical
        {
            std::cout << "Thread " << id
                      << " of " << count << '\n';
        }
    }
    std::cout << "Region finished\n";
}
\end{lstlisting}

\code{num_threads(4)} jest żądaniem liczby wątków, a faktyczna liczba
może zależeć od ograniczeń środowiska wykonawczego. \code{critical}
chroni tu przed przeplataniem fragmentów wypisywanego tekstu,
nie narzuca kolejności identyfikatorów wątków.

\subsection{Równoległa pętla \texttt{parallel for}}

Najczęstsze zastosowanie OpenMP polega na podziale \emph{niezależnych}
iteracji jednej pętli między uczestników zespołu.

\begin{lstlisting}[caption={Wypełnianie wektora: niezależne iteracje.}]
#include <cstddef>
#include <vector>

int main() {
    constexpr long long n = 100000;
    std::vector<double> squares(static_cast<std::size_t>(n));

    #pragma omp parallel for default(none) shared(squares, n) \
        schedule(static)
    for (long long i = 0; i < n; ++i) {
        squares[static_cast<std::size_t>(i)] =
            static_cast<double>(i) * i;
    }
}
\end{lstlisting}

\code{default(none)} wymaga jawnego określenia udziału zmiennych,
które nie mają z góry ustalonych atrybutów współdzielenia. Jest
wartościowym narzędziem kontroli poprawności. Zmienna sterująca
pętlą \code{i} jest prywatna na mocy reguł dyrektywy \code{for}.
Każdy wątek zapisuje inny element wektora; nie ma konfliktu,
zakładając że nikt nie zmienia równolegle rozmiaru ani struktury wektora.

\subsection{\texttt{shared}, \texttt{private}, \texttt{firstprivate}, \texttt{lastprivate}}

\begin{description}[style=nextline]
\item[\code{shared(x)}] Wątki odnoszą się do tego samego obiektu.
Dostępy konfliktowe wymagają ochrony lub podziału danych.
\item[\code{private(x)}] Każdy uczestnik ma własny egzemplarz
(z zachowaniem reguł inicjalizacji odpowiednich typów);
nie należy zakładać skopiowania początkowej wartości zmiennej pierwotnej.
\item[\code{firstprivate(x)}] Nowy egzemplarz zostaje zainicjalizowany
wartością odpowiadającej mu zmiennej sprzed regionu.
\item[\code{lastprivate(x)}] Po regionie wynik z logicznie ostatniej
iteracji (niekoniecznie wykonanej na końcu w czasie rzeczywistym)
zostaje przypisany do zmiennej pierwotnej.
\end{description}

\begin{lstlisting}[caption={Przykład firstprivate i lastprivate w OpenMP.}]
#include <iostream>

int main() {
    int seed = 10;
    int last = -1;

    #pragma omp parallel for firstprivate(seed) lastprivate(last)
    for (int i = 0; i < 8; ++i) {
        int local = seed + i;
        last = local;
    }
    std::cout << last << '\n'; // 17
}
\end{lstlisting}

Zmienna \code{seed} jest kopiowana do prywatnych egzemplarzy.
Wartość \code{last} po zakończeniu odpowiada iteracji $i=7$,
nie zaś temu wątkowi, który jako ostatni fizycznie zakończył pracę.

\subsection{Wyścig w redukcji i klauzula \texttt{reduction}}

Poniższa konstrukcja jest \textbf{błędna}:
\begin{lstlisting}[caption={NIEPOPRAWNIE: konfliktowe modyfikacje sumy.}]
// Fragment celowo bledny:
// double sum = 0;
// #pragma omp parallel for
// for (int i = 0; i < n; ++i) sum += values[i];
\end{lstlisting}

Każdy wątek modyfikuje wspólny \code{sum}. Możemy chronić dodawanie
blokadą, lecz w klasycznym sumowaniu wygodniejsze są prywatne sumy,
łączone pod koniec obliczeń:

\begin{lstlisting}[caption={Redukcja równoległa w OpenMP.}]
#include <iostream>
#include <vector>

int main() {
    std::vector<double> values(1000000, 1.0);
    const long long n = static_cast<long long>(values.size());
    double sum = 0.0;

    #pragma omp parallel for reduction(+:sum) schedule(static)
    for (long long i = 0; i < n; ++i) {
        sum += values[static_cast<std::size_t>(i)];
    }
    std::cout << sum << '\n'; // 1000000
}
\end{lstlisting}

OpenMP tworzy prywatne wartości redukcji i łączy je operatorem
\code{+}. Dostępne są również inne redukcje, np. \code{*},
\code{min} i \code{max} (zależnie od typu oraz konstrukcji).
Ważne: dla liczb zmiennoprzecinkowych kolejność działań może być
inna przy różnych zespołach i harmonogramach. Wyniki mogą zatem
różnić się na ostatnich cyfrach, bo
\[
 (a+b)+c\ne a+(b+c)
\]
może zachodzić w arytmetyce IEEE~754. Dla algorytmów numerycznych
mających wymagania reprodukowalności nie wystarczy zakładać
identycznych bitowo wyników każdej redukcji.

\subsection{Harmonogramowanie iteracji: \texttt{schedule}}

\begin{center}
\small
\begin{tabularx}{\textwidth}{@{}p{0.29\textwidth}X@{}}
\toprule
\textbf{Tryb} & \textbf{Znaczenie} \\
\midrule
\code{static} & Przydział iteracji ustalany na początku; mały narzut, dobry dla równych kosztów. \\
\code{static,k} & Przydział kolejnych porcji po $k$ iteracji w sposób cykliczny. \\
\code{dynamic,k} & Porcje po $k$ iteracji przydzielane pracownikom w trakcie pracy; dobre dla nierównych obciążeń, większy narzut. \\
\code{guided,k} & Rozmiary przydzielanych porcji zasadniczo maleją, do narzuconego minimum; kompromis między narzutem a równoważeniem. \\
\code{runtime} & Ustawienia pobierane z środowiska, m.in. \code{OMP_SCHEDULE}. \\
\bottomrule
\end{tabularx}
\end{center}

\begin{lstlisting}[caption={Równoważenie kosztów obliczeń w pętli.}]
#pragma omp parallel for schedule(dynamic, 8)
for (int i = 0; i < n; ++i) {
    perform_variable_cost_work(i);
}
\end{lstlisting}

Przy równomiernych obliczeniach \code{static} bywa tańszy od
\code{dynamic}. Dla zadań nieregularnych dynamiczny harmonogram
zmniejsza bezczynność pracowników, ale może pogorszyć lokalność pamięci.
Doboru dokonujemy na podstawie pomiarów, nie wyłącznie intuicji.

\subsection{Synchronizacja: \texttt{critical}, \texttt{atomic}, \texttt{barrier}, \texttt{nowait}}

\begin{description}[style=nextline]
\item[\code{critical}] Chroni region wykonywany jednocześnie najwyżej
przez jeden wątek w danej grupie nazwanych regionów krytycznych.
\item[\code{atomic}] Zapewnia atomową obsługę wspieranego wyrażenia
(np. inkrementacji); nie obejmuje dowolnego bloku instrukcji.
\item[\code{barrier}] Wymaga, żeby uczestnicy zespołu dotarli do
punktu synchronizacji, zanim którykolwiek przejdzie dalej.
\item[\code{nowait}] Usuwa domyślną barierę tam, gdzie klauzula jest
dopuszczalna; wymaga osobnego rozumowania o zależnościach danych.
\end{description}

\begin{lstlisting}[caption={Dwie alternatywy zwiększania licznika.}]
int counter = 0;

#pragma omp parallel for
for (int i = 0; i < 10000; ++i) {
    #pragma omp atomic update
    ++counter;
}
// counter == 10000
\end{lstlisting}

Dla tej operacji \code{reduction(+:counter)} zwykle ogranicza
rywalizację lepiej niż atomowa inkrementacja przy każdej iteracji.
Z kolei \code{critical} jest przydatny, gdy atomowo musi wykonać się
cały fragment kodu, a nie tylko jedno proste wyrażenie.

\begin{lstlisting}[caption={Bariera wewnątrz istniejącego zespołu.}]
#pragma omp parallel
{
    prepare_local_data();
    #pragma omp barrier
    consume_results_after_everyone_prepared();
}
\end{lstlisting}

Nie można wstawiać bariery w taki sposób, aby część wątków zespołu
ją pomijała, a pozostałe na niej czekały. Podobnie \code{nowait}
nie może być dodawane mechanicznie, jeśli następna pętla potrzebuje
wyników pierwszej.

\subsection{\texttt{sections} i \texttt{single}}

Gdy mamy kilka niezależnych, różnych funkcji, wygodne bywa
\code{sections}; gdy chcemy pojedynczego producenta zadań,
wybieramy \code{single}:

\begin{lstlisting}[caption={Kilka różnych zadań i jedna sekcja inicjalizacyjna.}]
#pragma omp parallel sections
{
    #pragma omp section
    load_input();

    #pragma omp section
    preprocess();
}

#pragma omp parallel
{
    #pragma omp single
    initialize_once();
    // Wszystkie watki mijaja koniec single po barierze.
}
\end{lstlisting}

Powyższy przykład \code{sections} ma sens tylko wtedy, gdy
\code{load_input()} i \code{preprocess()} są rzeczywiście niezależne.
W przeciwnym razie potrzebna jest synchronizacja albo zadania
z opisanymi zależnościami.

\subsection{Zadania OpenMP: \texttt{task}, \texttt{taskwait}, \texttt{taskgroup}}

Podział pętli \code{for} pasuje do regularnych obliczeń. Dla pracy
nieregularnej, np. rekursywnego przetwarzania drzewa, używa się
\code{task}. Dyrektywa tworzy jednostkę pracy, którą środowisko
wykonawcze może przydzielić pracownikowi zespołu. Nie jest to
równoznaczne z utworzeniem nowego wątku systemowego.

\begin{lstlisting}[caption={Tworzenie i oczekiwanie na zadania OpenMP.}]
#include <iostream>

int main() {
    int results[4]{};
    #pragma omp parallel
    {
        #pragma omp single
        {
            for (int i = 0; i < 4; ++i) {
                #pragma omp task firstprivate(i) shared(results)
                {
                    results[i] = i * i;
                }
            }
            #pragma omp taskwait
        }
    }
    for (int x : results) std::cout << x << ' ';
}
\end{lstlisting}

\code{taskwait} oczekuje na odpowiednie bezpośrednie zadania potomne
bieżącego zadania. Gdy chcemy objąć również ich potomków,
praktyczną konstrukcją jest \code{taskgroup}. Koniec regionu
\code{parallel} również zapewnia zakończenie prac należących
do tej struktury zadaniowej.

\subsection{Zależności pomiędzy zadaniami: \texttt{depend}}

\begin{lstlisting}[caption={Zadanie drugie korzysta z wyniku pierwszego.}]
#include <cassert>

int main() {
    int a = 0, b = 0;
    #pragma omp parallel
    {
        #pragma omp single
        {
            #pragma omp task shared(a) depend(out:a)
            {
                a = 21;
            }
            #pragma omp task shared(a,b) depend(in:a)
            {
                b = a * 2;
            }
            #pragma omp taskwait
        }
    }
    assert(b == 42);
}
\end{lstlisting}

Klauzule \code{depend(in:...)}, \code{depend(out:...)} i
\code{depend(inout:...)} opisują zależności danych m.in.
pomiędzy zadaniami rodzeństwa. Nie zastępują synchronizacji
przy dowolnych dostępach do innych, nieopisanych zasobów.

\subsection{\texttt{collapse} i \texttt{simd}: pętle wielowymiarowe i wektoryzacja}

\code{collapse(2)} łączy przestrzeń iteracji dwóch odpowiednio
zagnieżdżonych pętli, zwiększając liczbę możliwych porcji pracy.
\code{simd} natomiast przekazuje kompilatorowi informację,
że iteracje dają się wykonać w modelu SIMD. Wektoryzacja i
równoległość wielowątkowa są różnymi poziomami przyspieszania.

\begin{lstlisting}[caption={Tablica 2D: parallel for z collapse(2).}]
#include <vector>

int main() {
    const int rows = 400, cols = 400;
    std::vector<double> a(rows * cols);
    #pragma omp parallel for collapse(2) schedule(static)
    for (int i = 0; i < rows; ++i)
        for (int j = 0; j < cols; ++j)
            a[i * cols + j] = i + 0.5 * j;
}
\end{lstlisting}

\begin{lstlisting}[caption={Wektoryzacja bez obowiązkowego tworzenia wielu wątków.}]
void scale(double* dest, const double* src, int n, double alpha) {
    #pragma omp simd
    for (int i = 0; i < n; ++i)
        dest[i] = alpha * src[i];
}
\end{lstlisting}

Prawidłowość \code{simd} wymaga faktycznej niezależności iteracji,
w tym braku niedozwolonego nakładania się zapisów i odczytów.
Nie należy wprowadzać dyrektywy w kodzie o nieustalonym aliasingu
bez wcześniejszej analizy zależności.

\subsection{Zmienne środowiskowe OpenMP i diagnostyka}

\begin{center}
\small
\begin{tabularx}{\textwidth}{@{}p{0.3\textwidth}X@{}}
\toprule
\textbf{Zmienna} & \textbf{Rola} \\
\midrule
\code{OMP_NUM_THREADS} & Żądana liczba wątków w zespołach. \\
\code{OMP_DYNAMIC} & Pozwala środowisku dynamicznie dostosowywać liczbę wątków. \\
\code{OMP_SCHEDULE} & Harmonogram stosowany przy \code{schedule(runtime)}. \\
\code{OMP_PROC_BIND} & Strategia wiązania wątków z zasobami obliczeniowymi. \\
\code{OMP_PLACES} & Definiuje miejsca wykonania dla mechanizmu wiązania. \\
\code{OMP_MAX_ACTIVE_LEVELS} & Ogranicza liczbę aktywnych poziomów zagnieżdżonych regionów. \\
\bottomrule
\end{tabularx}
\end{center}

Przykładowa sesja uruchomieniowa:
\begin{lstlisting}[language=bash,caption={Sterowanie środowiskiem bez zmiany źródeł.}]
OMP_NUM_THREADS=8 OMP_SCHEDULE=dynamic,32 ./example
OMP_NUM_THREADS=4 OMP_PROC_BIND=close ./example
\end{lstlisting}

Rzeczywista dostępność i interpretacja opcji dotyczących wiązania
wątków zależą od środowiska wykonawczego i maszyny. Zagnieżdżone
regiony OpenMP oraz łączenie OpenMP z \code{std::thread}
może spowodować \emph{oversubscription}: więcej aktywnych
wątków roboczych niż zasobów obliczeniowych.

\subsection{OpenMP a model pamięci C++}

OpenMP określa własny model pamięci, operację \emph{flush} oraz
semantykę synchronizacji dyrektyw, takich jak \code{barrier},
\code{critical} i \code{atomic}. \emph{Flush} nie jest funkcją
\code{std::atomic_thread_fence}, a samo słowo ,,flush'' nie znaczy
,,wyczyść wszystkie cache procesora''.

W praktyce warto trzymać się jednej, konsekwentnej reguły:
\begin{itemize}
\item dane dzielone między iteracje OpenMP synchronizujemy przede
wszystkim przez odpowiednie konstrukcje OpenMP;
\item dane wymieniane pomiędzy własnymi wątkami C++ synchronizujemy
przez \code{std::mutex}, \code{std::atomic} lub inne narzędzia C++;
\item przy \emph{mieszaniu} obu modeli wymagane jest szczególnie
staranne rozumowanie o zgodności mechanizmów, zamiast zakładania,
że przypadkowe \code{flush} naprawi wyścig danych.
\end{itemize}

Wyjątków C++ nie należy wypuszczać poza regiony OpenMP w sposób
naruszający zasady specyfikacji. Jeżeli zadanie może zgłosić wyjątek,
trzeba obsłużyć go wewnątrz odpowiedniego regionu i przekazać
informację o błędzie w kontrolowany sposób po zakończeniu pracy.

\subsection{OpenMP czy \texttt{std::thread}, \texttt{async}, a może Ranges?}

\begin{center}
\small
\begin{tabularx}{\textwidth}{@{}p{0.28\textwidth}X@{}}
\toprule
\textbf{Mechanizm} & \textbf{Typowy przypadek użycia} \\
\midrule
\code{std::thread}/\code{jthread} & Długowieczny wątek, własny protokół komunikacji, kontrola czasu życia. \\
\code{std::async} & Niewiele niezależnych obliczeń z wynikiem w \code{future}. \\
Pula wątków & Wiele krótkich lub średnich zadań, ograniczona liczba pracowników. \\
OpenMP \code{parallel for} & Obliczenia naukowe i numeryczne z niezależnymi iteracjami. \\
OpenMP \code{task} & Dynamiczne drzewo prac i zależności. \\
\code{std::execution::par} & Algorytmy standardowe C++17+ z odpowiednią polityką wykonania i dostępnym backendem. \\
\bottomrule
\end{tabularx}
\end{center}

Polityka \code{std::execution::par} nie stanowi gwarancji faktycznego
wielowątkowego wykonania na każdym kompilatorze i bibliotece
standardowej. \emph{Ranges views} z C++20 są zazwyczaj przede
wszystkim leniwymi adaptatorami zakresów, nie automatycznym
mechanizmem równoległego planowania.

\subsection{Studium numeryczne: porównanie trzech sposobów sumowania}

Rozważmy trzy warianty: zwykły \code{std::accumulate}, równoległe
sumowanie przez \code{std::async} oraz OpenMP \code{reduction}.
W przykładzie\footnote{Przy dzieleniu dużego wektora na wiele części
należy uwzględnić balans obciążenia, rozmiary porcji i koszt
agregacji wyników.} obliczamy sumę kwadratów wejściowych liczb
przez OpenMP, a następnie porównujemy z referencyjną pętlą:

\begin{lstlisting}[caption={Praktyczny przykład: suma kwadratów OpenMP z pomiarem.}]
#include <cmath>
#include <iostream>
#include <omp.h>
#include <vector>

int main() {
    constexpr int n = 1000000;
    std::vector<double> a(n);
    for (int i = 0; i < n; ++i) a[i] = i * 1e-5;

    double sequential = 0.0;
    double t0 = omp_get_wtime();
    for (double x : a) sequential += x * x;
    double t1 = omp_get_wtime();

    double parallel = 0.0;
    #pragma omp parallel for reduction(+:parallel) schedule(static)
    for (int i = 0; i < n; ++i)
        parallel += a[i] * a[i];
    double t2 = omp_get_wtime();

    std::cout << "sequential: " << sequential
              << ", time: " << t1 - t0 << " s\n";
    std::cout << "OpenMP: " << parallel
              << ", time: " << t2 - t1 << " s\n";
    std::cout << "relative error: "
              << std::abs(sequential - parallel) /
                 (1.0 + std::abs(sequential)) << '\n';
}
\end{lstlisting}

\nota{Ten program jest demonstracją interfejsu, a nie metodologicznie
pełnym benchmarkiem. Należy wykonać rozgrzewkę, wiele powtórzeń,
uwzględnić rozrzut wyników i sprawdzić, czy kompilator nie usunął
obliczeń. Koszt dostępu do pamięci może ograniczać przyspieszenie;
przy małym \code{n} równoległa wersja może być wolniejsza.}



\section{Porównanie i zasady projektowania}
\label{sec:choose}

Po zapoznaniu się z modelem pamięci i OpenMP potrafimy rozróżnić
trzy poziomy pracy: \emph{własność wątku} (\code{thread}/\code{jthread}),
\emph{własność zadania i jego wyniku} (\code{future}, pule) oraz
\emph{deklaratywny podział iteracji/zadań} (OpenMP). Wybór nie
powinien zależeć jedynie od liczby linijek kodu.


\begin{center}
\small
\begin{tabularx}{\textwidth}{@{}lX@{}}
\toprule
\textbf{Narzędzie} & \textbf{Kiedy używać} \\
\midrule
\code{thread} & Potrzebna jest bezpośrednia kontrola nad wątkiem i jego czasem życia. \\
\code{jthread} & W C++20, gdy przydaje się RAII i kooperacyjny sygnał zatrzymania. \\
\code{mutex} & Kilka operacji musi spójnie chronić współdzielony stan. \\
\code{condition_variable} & Wątek powinien zasnąć do spełnienia warunku na stanie. \\
\code{atomic} & Niezależne, atomowe operacje na zmiennej lub starannie zaprojektowana synchronizacja. \\
\code{promise} & Producent sam decyduje, kiedy dostarczyć wynik albo wyjątek. \\
\code{future} & Odbiór jednego wyniku i ewentualne oczekiwanie. \\
\code{shared_future} & Kilku odbiorców potrzebuje tej samej wartości. \\
\code{async} & Chcemy uruchomić funkcję i odebrać jej wynik bez zarządzania wątkiem. \\
\code{packaged_task} & Trzeba oddzielić zadanie od mechanizmu jego uruchamiania. \\
Pula wątków & Dużo zadań i ograniczona liczba pracowników; kontrolowane zamykanie. \\
OpenMP & Duże obliczenia z niezależnymi iteracjami albo grafem zadań. \\
\bottomrule
\end{tabularx}
\end{center}

\subsection{Najczęstsze błędy}

\begin{enumerate}
\item Użycie zwykłej zmiennej współdzielonej bez synchronizacji
(\emph{data race}, niezdefiniowane zachowanie).
\item Założenie, że \code{++x} na zwykłym \code{int} jest atomowe.
\item Zniszczenie aktywnego \code{std::thread} bez \code{join()} lub \code{detach()}.
\item Przekazanie referencji do obiektu, którego czas życia kończy się
przed zakończeniem pracy wątku.
\item Oczekiwanie przez \code{condition_variable} bez predykatu.
\item Mylenie powiadomienia z trwałym zapisem informacji o stanie.
\item Założenie, że domyślne \code{std::async(f)} zawsze uruchamia nowy wątek.
\item Natychmiastowe niszczenie tymczasowego \code{future} z \code{std::async}.
\item Traktowanie \code{wait_for()} jako anulowania zadania.
\item Próba wielokrotnego \code{get()} na zwykłym \code{future}.
\item Nadmierne używanie \code{memory_order_relaxed} bez dowodu poprawności.
\item Zakładanie, że dla każdego drobnego zadania opłaca się uruchomić
osobny wątek.
\item Przenoszenie \code{memory_order_relaxed} do algorytmu publikującego dane
bez analizy zależności między wątkami.
\item Stosowanie \code{\#pragma omp parallel for} do pętli z zależnościami
między iteracjami albo wspólnym niechronionym akumulatorem.
\item Mieszanie własnej puli wątków z OpenMP bez kontroli liczby aktywnych
pracowników; pojawia się ryzyko \emph{oversubscription}.
\end{enumerate}

\subsection{Standard języka a dostępność API}

\begin{center}
\begin{tabularx}{\textwidth}{@{}lX@{}}
\toprule
\textbf{Standard} & \textbf{Wybrane mechanizmy} \\
\midrule
C++11 & \code{thread}, \code{mutex}, \code{condition_variable}, \code{atomic}, \code{promise}, \code{future}, \code{async}, \code{packaged_task}, \code{call_once}, \code{thread_local}. \\
C++14 & M.in. ułatwienia związane z lambdami, np. inicjalizujące przechwycenie użyte w przykładzie \code{promise}. \\
C++17 & \code{scoped_lock}, \code{shared_mutex}, \code{shared_lock} (ta ostatnia od C++14). \\
C++20 & \code{jthread}, \code{stop_token}, \code{atomic::wait}, \code{semaphore}, \code{latch}, \code{barrier}, \code{osyncstream}. \\
\bottomrule
\end{tabularx}
\end{center}

Przykład z inicjalizującym przechwyceniem:
\begin{center}
\code{[p = std::move(producer)]}
\end{center}
wymaga C++14, mimo że sam \code{std::promise} istnieje od C++11.

\section{Kompilacja, sprawdzanie i diagnostyka}

Każdy kompletny przykład można zapisać w osobnym pliku z rozszerzeniem
\code{.cpp}. Dla kodu C++20 typowe wywołanie kompilatora GCC lub Clang to:

\begin{lstlisting}[language=bash,caption={Kompilacja przykładu i uruchomienie programu.}]
g++ -std=c++20 -O2 -Wall -Wextra -pthread example.cpp -o example
./example
\end{lstlisting}

Podczas nauki szczególnie pomocny jest \emph{ThreadSanitizer}, który
może wykrywać wiele wyścigów danych (dostępność zależy od platformy
oraz kompilatora):

\begin{lstlisting}[language=bash,caption={Przykładowe uruchomienie z ThreadSanitizer.}]
g++ -std=c++20 -O1 -g -pthread -fsanitize=thread example.cpp -o example_tsan
./example_tsan
\end{lstlisting}

Nawet jeśli test przejdzie tysiąc razy, nie dowodzi to braku błędów
współbieżności. Poprawność wymaga analizy synchronizacji i czasu życia;
narzędzia diagnostyczne stanowią jej uzupełnienie.

\section*{Podsumowanie rozdziału}
\addcontentsline{toc}{section}{Podsumowanie rozdziału}

Z \code{std::thread} i \code{std::jthread} sterujemy \emph{wątkami},
a z \code{std::async} i \code{std::packaged_task} opisujemy raczej
\emph{zadania}. \code{future}/\code{promise} tworzą interfejs dostarczania
i odbioru jednego wyniku. \code{mutex} i \code{condition_variable}
pozwalają chronić współdzielony stan i czekać na jego zmianę.
\code{atomic} udostępnia operacje atomowe oraz mechanizmy porządkowania
pamięci. Algorytmy \emph{lock-free} wymagają ponadto zarządzania czasem
życia węzłów i analizy problemu ABA. Pula wątków separuje zadania
od fizycznych pracowników, a OpenMP udostępnia deklaratywny podział
pracy na iteracje, redukcje i zadania. Żadne z tych narzędzi nie
zastępuje pozostałych we wszystkich zastosowaniach.

Dobra architektura zaczyna się od pytania, \emph{które dane są współdzielone,
kto jest właścicielem zadania, kiedy wolno zakończyć jego pracę oraz
jak przekazać wynik i wyjątek}. Dopiero potem należy wybrać konkretne
narzędzia biblioteki standardowej.

\section*{Literatura i dokumentacja}
\addcontentsline{toc}{section}{Literatura i dokumentacja}

\begin{itemize}
\item \url{https://cppreference.com/cpp/thread/async} --- \code{std::async} i polityki uruchamiania.
\item \url{https://en.cppreference.com/w/cpp/thread/future} --- \code{std::future}.
\item \url{https://cppreference.com/cpp/thread/condition_variable/wait} --- oczekiwanie z predykatem.
\item \url{https://en.cppreference.com/w/cpp/atomic/atomic/wait} --- \code{atomic::wait} w C++20.
\item \url{https://en.cppreference.com/w/cpp/atomic/memory_order} --- porządek pamięci.
\item \url{https://en.cppreference.com/w/cpp/atomic/atomic/compare_exchange} --- CAS.
\item \url{https://www.openmp.org/specifications/} --- oficjalne wersje i dokumentacja OpenMP (w tym 5.2 i 6.0).
\item \url{https://www.openmp.org/spec-html/5.2/openmp.html} --- treść specyfikacji OpenMP 5.2.
\item \url{https://www.openmp.org/spec-html/5.2/openmpse4.html} --- model pamięci OpenMP.
\item Anthony Williams, \emph{C++ Concurrency in Action}, wyd. 2, Manning, 2019.
\end{itemize}
